\SPBGContentSlide
  {Постановка задачи исследования}
  {%
    \Body
    Дано множество объектов, каждый из которых описывается вектором признаков
    и меткой одного из двух классов:
    \[
      (p_1,\xi_1),\ldots,(p_m,\xi_m),
      \qquad
      p_j\in\mathbb{R}^n,\quad \xi_j\in\{-1,+1\}.
    \]

    Необходимо построить линейный классификатор
    \[
      \widehat{\xi}(p)=\sgn(\langle w,p\rangle+\beta),
    \]
\begin{itemize}
  \item \(p\) -- вектор признаков объекта;
  \item \(w\) -- вектор параметров классификатора;
  \item \(\beta\) -- свободный член, задающий смещение разделяющей гиперплоскости;
  \item \(\langle w,p\rangle+\beta=0\) -- разделяющая гиперплоскость;
  \item \(\sgn(\cdot)\) -- правило отнесения объекта к одному из двух классов.
\end{itemize}

  }